\documentclass[aspectratio=169,10pt,t]{beamer}
\usepackage{ac1}
\aulainfo{04}{Pensamento Algorítmico}{Algoritmos, refinamentos sucessivos, fluxogramas e pseudocódigo}

\begin{document}
\capa

\begin{frame}{Roteiro}
  \begin{itemize}
    \item Linguagem de programação e algoritmos (Cap.~2.1--2.2)
    \item Um algoritmo computacional simples: a busca no dicionário (2.3)
    \item Refinamentos sucessivos
    \item Fluxogramas (2.4)
    \item Pseudocódigo (2.5)
    \item \textbf{Prática:} algoritmos no papel
  \end{itemize}
\end{frame}

\begin{frame}{Algoritmo: passos não ambíguos}
  \begin{block}{O computador executa \emph{exatamente} o que você escreve}
    Ele não ``entende a intenção''. Cada passo precisa ser preciso e executável.
  \end{block}
  \begin{columns}[T]
    \column{0.48\textwidth}
    \textbf{Ambíguo}
    \begin{itemize}
      \item ``Misture bem''
      \item ``Espere um pouco''
      \item ``Ache a palavra''
    \end{itemize}
    \column{0.48\textwidth}
    \textbf{Preciso}
    \begin{itemize}
      \item ``Misture por 30 segundos''
      \item ``Espere 5 minutos''
      \item ``Abra o livro na página do meio e compare''
    \end{itemize}
  \end{columns}
\end{frame}

\begin{frame}{Um algoritmo computacional: achar uma palavra no dicionário}
  \textbf{Primeira tentativa (força bruta):}
  \begin{enumerate}
    \item Abra na primeira página
    \item Se a palavra estiver ali, pare
    \item Senão, vire a página e volte ao passo 2
  \end{enumerate}
  Funciona, mas é lento: um dicionário com 1000 páginas pode exigir \(1000\) passos.
  \vspace{1ex}
  \begin{alertblock}{Pergunta}
    Como uma pessoa \emph{de verdade} procura no dicionário? Isso explora o fato de ele estar \textbf{ordenado}.
  \end{alertblock}
\end{frame}

\begin{frame}{Refinamento: a busca binária}
  \begin{enumerate}
    \item Abra o dicionário no meio
    \item Se a palavra está na página, \alert{pare}
    \item Se a palavra vem \emph{antes} da página atual, descarte a metade de depois
    \item Se vem \emph{depois}, descarte a metade de antes
    \item Repita com a metade que sobrou
    \item Se não sobrar nenhuma página, a palavra \alert{não existe}
  \end{enumerate}
  \vspace{1ex}
  \begin{exampleblock}{Eficiência}
    Cada passo divide o problema por 2: com 1000 páginas bastam \(\lceil\log_2 1000\rceil = 10\) passos.
  \end{exampleblock}
\end{frame}

\begin{frame}{Refinamentos sucessivos}
  Estratégia para problemas grandes: começar com passos de \textbf{alto nível} e detalhar cada um.
  \vspace{1ex}
  \begin{columns}[T]
    \column{0.45\textwidth}
    \textbf{Nível 1}
    \begin{enumerate}
      \item Ler as notas
      \item Calcular a média
      \item Informar a situação
    \end{enumerate}
    \column{0.52\textwidth}
    \textbf{Nível 2 (refinando o passo 3)}
    \begin{enumerate}
      \item[3.1] Se média \(\ge 7\): ``aprovado''
      \item[3.2] Senão, se média \(\ge 4\): ``prova final''
      \item[3.3] Senão: ``reprovado''
    \end{enumerate}
  \end{columns}
  \vspace{1ex}
  \footnotesize Só quando todos os passos são simples o suficiente é que escrevemos o código.
\end{frame}

\begin{frame}{Fluxogramas: símbolos básicos}
  \begin{center}
  \begin{tikzpicture}[
      term/.style={draw=uerjazul,rounded corners=8pt,minimum width=2.6cm,minimum height=0.8cm,fill=uerjcinza,font=\small},
      proc/.style={draw=uerjazul,minimum width=2.6cm,minimum height=0.8cm,fill=uerjcinza,font=\small},
      io/.style={draw=uerjazul,trapezium,trapezium left angle=70,trapezium right angle=110,minimum width=2.6cm,minimum height=0.8cm,fill=uerjcinza,font=\small},
      dec/.style={draw=uerjazul,diamond,aspect=2,fill=uerjcinza,font=\small}]
    \node[term] (a) at (0,0) {Início / Fim};
    \node[proc] (b) at (4,0) {Processamento};
    \node[io]   (c) at (8,0) {Entrada / Saída};
    \node[dec]  (d) at (12,0) {Decisão?};
  \end{tikzpicture}
  \end{center}
  \begin{itemize}
    \item Setas indicam o \textbf{fluxo} de execução
    \item Decisão tem duas saídas: \emph{sim} e \emph{não}
    \item Um fluxograma bom cabe em uma página e se lê de cima para baixo
  \end{itemize}
\end{frame}

\begin{frame}{Exemplo: situação do aluno}
  \begin{center}
  \begin{tikzpicture}[scale=0.68,transform shape,node distance=0.7cm,
      term/.style={draw=uerjazul,rounded corners=8pt,minimum width=2.4cm,minimum height=0.6cm,fill=uerjcinza,font=\footnotesize},
      proc/.style={draw=uerjazul,minimum width=3.3cm,minimum height=0.6cm,fill=uerjcinza,font=\footnotesize},
      dec/.style={draw=uerjazul,diamond,aspect=3,fill=uerjcinza,font=\footnotesize,inner sep=1pt},
      seta/.style={-{Stealth},thick,uerjazul}]
    \node[term] (ini) {Início};
    \node[proc,below=of ini] (le) {Ler nota1 e nota2};
    \node[proc,below=of le] (m) {media $\leftarrow$ (nota1 + nota2)/2};
    \node[dec,below=of m] (d1) {media $\ge$ 7?};
    \node[proc,right=2.2cm of d1] (ap) {``Aprovado''};
    \node[dec,below=of d1] (d2) {media $\ge$ 4?};
    \node[proc,right=2.2cm of d2] (pf) {``Prova final''};
    \node[proc,below=of d2] (rp) {``Reprovado''};
    \node[term,right=2.2cm of rp] (fim) {Fim};
    \draw[seta] (ini)--(le); \draw[seta] (le)--(m); \draw[seta] (m)--(d1);
    \draw[seta] (d1)--node[above,font=\scriptsize]{sim}(ap);
    \draw[seta] (d1)--node[left,font=\scriptsize]{não}(d2);
    \draw[seta] (d2)--node[above,font=\scriptsize]{sim}(pf);
    \draw[seta] (d2)--node[left,font=\scriptsize]{não}(rp);
    \draw[seta] (rp)--(fim);
    \draw[seta] (ap) -- ++(2.3,0) |- (fim);
    \draw[seta] (pf) -- ++(2.3,0) |- (fim);
  \end{tikzpicture}
  \end{center}
\end{frame}

\begin{frame}{Pseudocódigo}
  Descrição \textbf{textual e estruturada} do algoritmo, sem a sintaxe rígida de uma linguagem.
  \vspace{1ex}
  \begin{columns}[T]
    \column{0.52\textwidth}
    \begin{block}{Pseudocódigo}
      \ttfamily\footnotesize
      LEIA nota1, nota2\\
      media $\leftarrow$ (nota1 + nota2) / 2\\
      SE media $\ge$ 7 ENTÃO\\
      \hspace*{1em}ESCREVA ``Aprovado''\\
      SENÃO SE media $\ge$ 4 ENTÃO\\
      \hspace*{1em}ESCREVA ``Prova final''\\
      SENÃO\\
      \hspace*{1em}ESCREVA ``Reprovado''
    \end{block}
    \column{0.44\textwidth}
    \begin{itemize}
      \item Foco na lógica, não na sintaxe
      \item Cada comando é uma ação clara
      \item Indentação mostra a estrutura
      \item Depois é traduzido para Python quase linha a linha
    \end{itemize}
  \end{columns}
\end{frame}

\begin{frame}[fragile]{Do pseudocódigo ao Python}
\begin{lstlisting}
nota1 = float(input("Primeira nota: "))
nota2 = float(input("Segunda nota: "))
media = (nota1 + nota2) / 2
if media >= 7:
    print("Aprovado")
elif media >= 4:
    print("Prova final")
else:
    print("Reprovado")
\end{lstlisting}
  \begin{itemize}
    \item \texttt{input} lê texto; \texttt{float(\ldots)} converte para número
    \item \texttt{elif} = ``senão se''
  \end{itemize}
  \fonte{code/Cap3/prog.3.3\_calcula\_final\_reprova.py}
\end{frame}

\begin{frame}{Prática: uma máquina só de somar e subtrair}
  Imagine uma máquina que só sabe \textbf{somar} e \textbf{subtrair}, e comparar com zero.

  \begin{exampleblock}{Desafio (Exercícios 2.4 a 2.6)}
    Escreva, em pseudocódigo, um algoritmo para
    \begin{enumerate}
      \item multiplicar \(a \times b\)
      \item calcular \(a^b\)
      \item dividir \(a \div b\) (quociente e resto)
    \end{enumerate}
  \end{exampleblock}
  \textbf{Dica para a multiplicação:} \(a\times b = a + a + \cdots + a\) (\(b\) vezes). Use um contador e um acumulador.
  \vspace{1ex}

  \footnotesize Depois desenhe os fluxogramas (Exercício 2.7).
\end{frame}

\begin{frame}[fragile]{Solução esperada: multiplicação por somas}
\begin{lstlisting}
a = int(input("a = "))
b = int(input("b = "))
produto = 0
contador = 0
while contador < b:
    produto = produto + a
    contador = contador + 1
print(a, "x", b, "=", produto)
\end{lstlisting}
  \begin{itemize}
    \item Já usamos uma repetição (\texttt{while}); ela será estudada na Aula 8
    \item Teste na mão: \(a=3\), \(b=4\) gera 0, 3, 6, 9, 12
  \end{itemize}
\end{frame}

\exercicios{2}{%
  \begin{itemize}\small
    \item Algoritmos do dia a dia: \textbf{2.1, 2.2}
    \item Busca binária e máquina de somar: \textbf{2.3 a 2.6}
    \item Fluxogramas: \textbf{2.7}
    \item Busca binária mais enxuta: \textbf{2.19}
  \end{itemize}}

\resumo{
  \item Algoritmo: passos finitos, precisos e sem ambiguidade
  \item Refinamentos sucessivos: de passos gerais a passos simples
  \item Fluxograma (visual) e pseudocódigo (textual) servem para pensar \emph{antes} de programar
  \item Próxima aula: identificadores, dados e tipos em Python
}

\end{document}
